\section{The heat field and a counted pair}
\label{sec:framework}

The local interaction of a pair of zeros is determined by the entire
exterior divisor, including nonreal zeros and their multiplicities.
The useful coordinates for this interaction are the center $c$ and the
signed discriminant $D$, rather than individually labelled target roots.
Our starting point is the pair law established in the base manuscript
and its cluster-transport development \cite{BakerBase,BakerGT41}.
The present paper extends its rational moments to every integer order,
relates them to split-pole moments, and then treats a complex exponent.
The Gaussian-resolvent predecessor \cite{BakerGR02,BakerGR05,BakerGR07}
supplies the first higher moment law and the quantitative comparison
used in Section~\ref{sec:bounds}.

The distinction between an identity and an estimate is essential to the
development. An exact representation of the exterior interaction is
valuable for organizing its evolution. A better gap estimate additionally
requires a quantitatively useful signed estimate for the relevant
coupled quantity, on a time interval whose boundary count and separation
have been established independently.
The later all-real comparison uses an additional source of sign: each
exterior denominator of an adjacent real pair is positive. This gives
an explicit local backward allowance from neighboring gaps, while its
statistical application still requires information about rare exceptions.

The development begins with exact structure. Sections~\ref{sec:integer}
and \ref{sec:split} derive the integer hierarchy and its split-pole
coordinates. Section~\ref{sec:lattice} places those rows on the same zero
lattice and proves their overlap and block-flux identities. This gives the
compatibility question in Section~\ref{sec:finite} a concrete spatial
meaning. Complex exponents, contours and signed forcing then lead to the
joint dynamics in Sections~\ref{sec:bounds}--\ref{sec:joint}.

The second stage concerns quantitative use. The backward and regional
laws of Section~\ref{sec:statistics} are followed by a joint-in-time
barrier in Section~\ref{sec:jointtime}. Section~\ref{sec:canonicaltime}
compares the initial-data and joint consumers on the same canonical
positive-time source, retaining all height and time errors. The compact
contour method in Section~\ref{sec:compact} separates source acquisition
from a noncircular stopping argument. Section~\ref{sec:forward-statistics}
then asks what finite time-zero geometry can acquire in the forward
direction. Section~\ref{sec:retained-forcing} retains the physical forcing
in the backward comparison and computes the exact cost of an unchanged
forward--backward excursion. These results motivate the common-source
two-sided refinement in Section~\ref{sec:two-sided}: a discrete enclosure
construction and a continuous residual identity, with their unacquired
strict-gain requirement stated explicitly. Supporting certificate-based
examples are deferred to Appendix~\ref{app:compact-example}; none of the
general theorems depends on those numerical values.

Strictly separating contours suffice for the local algebra, lattice and
compact stopping results. The quantitative moment-order, complex-power
and truncation estimates retain the stronger collar fixed below. Stating
these domains separately prevents an adjacent-pair or negative-time
application from silently importing a stronger geometric hypothesis.

\subsection{Canonical normalization and hypotheses}

For $u\ge0$, set
\begin{equation}
 r_k(u)=\pi k^2\e^{4u},\qquad
 \Phi(u)=\sum_{k\ge1}\e^u(2r_k(u)^2-3r_k(u))\e^{-r_k(u)},\qquad
 H(t,z)=\int_0^\infty\e^{tu^2}\Phi(u)\cos(zu)\dd u.
 \label{eq:source}
\end{equation}
This is the unchanged complete theta source, with the normalization of
\cite{BakerBase}; in particular $H(0,z)=\xi((1+\ii z)/2)/8$.
Its positive half-line source gives $H(t,0)>0$ for real $t$.
On compact sets in $(t,z)$ the integrand, including every fixed derivative,
is dominated by a constant times a polynomial in $u$ times
\[
 \exp\{9u+\Re(t)u^2+|\Im z|u-\pi\e^{4u}\}.
\]
Consequently $H$ is jointly entire, real on the real plane, even in $z$,
and satisfies the physical-time equation
\begin{equation}
 H_t=-H_{zz}.
 \label{eq:heat}
\end{equation}
The sign of time in this equation is fixed throughout the paper.
Maximizing $ru-(\pi/4)\e^{4u}$ after reserving an integrable part
of the source bound gives
$\log\max_{|z|\le r}|H(t,z)|=O(1+r\log(2+r))$
uniformly for $t$ in any fixed compact real interval. The bounded factor
$\Re(t)u^2$ is absorbed by the reserved double-exponential decay; at the
maximizing scale it contributes only $O((\log(2+r))^2)$. Thus the order
is at most one on every such interval. This does not extend any later
positive-time localization or numerical error constant to negative time.

For the quantitative strong-collar estimates,
fix a compact real interval $I=[t_0,t_1]\subset[0,1/2]$ on which an independently
established count-two tube isolates a conjugation-invariant pair. Freeze
the height and collar parameters
\begin{equation}
 \begin{gathered}
 T\ge64,\quad L=\log T,\quad \ell=L^{-1},\quad d\ge0,
 \quad \max(2,4\sqrt d)\le a\le L,\quad r_0=a\ell,\\
 c(t)\in[T,2T],\quad D(t)\in\R,\quad |D(t)|\le d\ell^2,\qquad
 z_\pm=c\pm s,\quad s^2=D,\\
 |z_\pm-c|\le r_0/4,\qquad
 w_j=\rho_j-c,\quad |w_j|\ge r_0
 \quad\text{for every exterior occurrence }\rho_j.
 \end{gathered}
 \label{eq:domain}
\end{equation}
An occurrence means a zero counted with its integer multiplicity.
The two target occurrences coincide when $D=0$; $D<0$ describes
a conjugate nonreal pair. The fixed scale $\ell$ is not differentiated.

We use the complete strip, unit-count and safe-line bounds inherited
from the canonical analysis in \cite{BakerBase,BakerGT41}. Specifically,
the base manuscript gives the projected count in ``Uniform counts,
including time zero'' and the safe-line estimate in ``A convex balance
with conjugate dissipation.'' With one
constant $\Cu\ge1$, enlarged to cover these inputs,
\begin{equation}
 \begin{gathered}
 |\Im\rho|\le1,\qquad
 \#\{\rho:v\le\Re\rho\le v+1\}\le\Cu\log(2+|v|),\\
 \left|\frac{H_z}{H}(t,v\pm4\ii)\right|
       \le\Cu\log(2+|v|),\qquad v\in\R,\quad t\in I.
 \end{gathered}
 \label{eq:count-safe}
\end{equation}
These are stated inputs, not conclusions of the moment identities.
In particular, those identities do not create or prolong the tube.

Write
\begin{equation}
 Q(t,w)=w^2-D(t),\quad Q_*(t,z)=(z-c(t))^2-D(t),\quad
 Q_j=Q(t,w_j),\quad \mu_j=\frac{\ell^2}{Q_j},\quad q_0=\frac{16}{15}.
 \label{eq:quadratic}
\end{equation}
The collar implies
\begin{equation}
 |Q_j|\ge\frac{15}{16}|w_j|^2,\qquad
 |\mu_j|\le\gamma:=\frac{q_0}{a^2}<1.
 \label{eq:collar-kernel}
\end{equation}
No right-half-plane condition is imposed on $Q_j$.

\subsection{Symmetric coordinates and the pair equation}

At a reference time $t_*$ choose the positive circle
$\Gamma=\{|z-c(t_*)|=r_0/2\}$, fixed while differentiating in $t$.
It has positive distance from all zeros. Its count remains two for
nearby complex times. If $B=H_z/H$, the argument principle gives
\begin{equation}
 J_k=\frac1{2\pi\ii}\oint_\Gamma z^kB(t,z)\dd z,
 \qquad c=J_1/2,\qquad D=J_2/2-c^2.
 \label{eq:coordinates}
\end{equation}
These formulas define holomorphic coordinates even at a double zero.
The factorization $H=Q_*A$ has $A$ jointly holomorphic and nonzero
inside $\Gamma$. This follows either from division by the counted
quadratic factor or from the Cauchy extension of the analytic boundary
quotient. Put $g=A_z/A$ there. Inserting this factorization into
\eqref{eq:heat}, with $w=z-c$, gives
\begin{equation}
 \frac{A_t}{A}=-(g_z+g^2)
       +\frac{2wc'+D'-2-4wg(t,c+w)}{w^2-D}.
 \label{eq:deflated}
\end{equation}
The numerator must vanish to the target multiplicities. Interpolating
$h_0(w)=wg(t,c+w)$ at $\pm s$, with Hermite interpolation at $s=0$,
therefore yields
\begin{equation}
 \begin{split}
 c'&=g(t,c-s)+g(t,c+s),\\
 D'&=2+4Dg[t;c-s,c+s]=2-4DM,\qquad
 M:=-g[t;c-s,c+s],\\
 \frac{A_t}{A}&=-(g_z+g^2)-4h_0[-s,s,w].
 \end{split}
 \label{eq:pairlaw}
\end{equation}
Divided differences are understood by analytic continuation at repeated
nodes. At a merger $M=-g_z(t,c)$ and $D'=2$. Thus a collision creates
no singularity in the signed discriminant equation.

\subsection{Rational traces and the complete divisor}

\begin{proposition}[Exterior traces]\label{prop:trace}
For $n\ge1$ and $k\ge2$, respectively, the moments
\begin{equation}
 S_n=\sum_j Q_j^{-n},\qquad T_k=\sum_j w_jQ_j^{-k}
 \label{eq:moments}
\end{equation}
are absolutely convergent and real on $I$. They have the local
holomorphic contour expressions
\begin{equation}
 S_n=-\frac1{2\pi\ii}\oint_\Gamma Q_*^{-n}B\dd z,
 \qquad
 T_k=-\frac1{2\pi\ii}\oint_\Gamma(z-c)Q_*^{-k}B\dd z.
 \label{eq:traces}
\end{equation}
In particular $M=S_1$. Their time derivatives exist through arbitrary
exterior mergers and a double target.
\end{proposition}

\begin{proof}
The unit count implies $\sum_\rho|\rho|^{-2}<\infty$.
The order bound and evenness give the canonical product
\begin{equation}
 H(t,z)=H(t,0)\prod_{\{\rho,-\rho\}}
              (1-z^2/\rho^2),
 \label{eq:product}
\end{equation}
where opposite groups, their conjugates and multiplicities are retained.
Indeed the genus-one factors of an opposite pair multiply to the displayed
factor, and the only possible linear exponential prefactor is constant
by evenness. The grouped logarithmic derivative converges locally away
from the divisor. The finite collar and the unit count give absolute
convergence of the moments; the slowest tails are $|\rho|^{-2}$
and $|\rho|^{-3}$.

For either kernel $R(w)=Q(w)^{-n}$ or $wQ(w)^{-k}$, the residue theorem
for the rational function $R(z-c)/(z-\rho_j)$ gives
\[
 \frac1{2\pi\ii}\oint_\Gamma
                 \frac{R(z-c)}{z-\rho_j}\dd z=-R(w_j).
\]
There is no residue at infinity. The genus-one additive constants
integrate to zero because the sum of the residues of $R$ is zero.
The same is true of the target contribution $R\,\partial_z Q_*/Q_*$:
it has poles only at the targets and decays faster than $1/z$.
Its combined target residue vanishes, even though the two residues
need not vanish individually. Uniform convergence on $\Gamma$ now
proves \eqref{eq:traces}. Taking $n=1$ also gives
$S_1=-g[t;c-s,c+s]$. The fixed-circle formulas and
\eqref{eq:coordinates} prove local analyticity without labelling
individual exterior roots. Reality follows from conjugation.
\end{proof}

\begin{lemma}[Passage from polynomial identities at a time slice]
\label{lem:slice-passage}
Suppose an identity for a fixed finite collection of the moments
\eqref{eq:moments}, $D$, and their first time derivatives is polynomial
in absolutely convergent moments, after all nonabsolute terms have
cancelled. If it holds for every finite polynomial heat field with an
isolated target pair, it holds for $H$ under \eqref{eq:domain}--
\eqref{eq:count-safe}, including multiple zeros.
\end{lemma}

\begin{proof}
First, for a finite polynomial at $t_*$, perturb its roots to distinct
ones preserving the count on $\Gamma$, and initialize its own polynomial
heat evolution there. The derivative
\begin{equation}
 B_t=-\frac{H_{zzz}}H+\frac{H_zH_{zz}}{H^2}
 \label{eq:log-time}
\end{equation}
and the fixed-circle coordinate and trace formulas depend continuously
on the polynomial and its first three spatial derivatives. The contour
denominators stay bounded away from zero. Thus an identity proved for
simple roots passes to all finite multiplicities.

Now truncate \eqref{eq:product} at increasing radii, retaining all ties,
opposite groups and conjugates. Include the targets for all sufficiently
large cutoffs. The product tail on a fixed compact set is
$O(\log N/N)$ by the unit count. The polynomials and every fixed
spatial derivative converge locally uniformly to $H(t_*,\cdot)$.
Give each truncated polynomial its own heat evolution initialized at
$t_*$. Formula \eqref{eq:log-time} and the trace formulas give convergence
of the values and first time derivatives of the stated observables at
that time. The remaining exterior sums converge absolutely. Passing
the identity to the limit proves it at $t_*$; repeat at every time
of $I$, with inward derivatives at its endpoints. This does not require
a common time interval for the truncated heat evolutions, and gives
no estimate uniform as the moment order tends to infinity.
\end{proof}
